\section*{Part X: Bayes Institute: A Governance System That Evolves With Risk}
\addcontentsline{toc}{section}{Part X: Bayes Institute: A Governance System That Evolves With Risk}

\subsection{61. Bayes's legal reality comes before its future title}

Bayes is an early venture concept, not currently a legally recognized university or established research campus. The founder has relevant technical expertise and a long academic ambition, but the company's exact current legal structure, shareholder records, tax registrations and employer-conflict permissions must be verified from documents. This chapter must not assume that the registered company and the Bayes brand are already legally aligned.

The 1st governance task is a \textbf{corporate facts audit}, not a university advisory board announcement.

\subsection{62. Phase G0: founder-led commercial compliance}

Verify legal identity, ownership, bank authority, tax registration, employment conflict, IP ownership, website claim accuracy, customer contracts, privacy/consent process, bookkeeping and statutory filings. Create a 13-week cash forecast, a single contract repository and a register of major decisions. Do not open a complex nonprofit or foreign subsidiary before evidence of demand requires it.

Keep the household, employer and venture finances strictly separate. Every customer deposit is accompanied by a delivery obligation.

\subsection{63. Phase G1: 1st paid programs}

Publish accurate terms and assessment criteria, require clear instructor and contractor agreements, obtain appropriate learner consents, establish refund and grievance routes, reconcile gateway collections, and create an independent 2nd-review path for consequential disputes. Assign an approval threshold for discounts, unusual promises and refunds. Track performance and consented proof honestly.

This is the small-company form of institutional legitimacy: buyers can understand, challenge and verify what they are purchasing.

\subsection{64. Phase G2: 1st college and enterprise partnerships}

Use defined statements of work, payment schedules, faculty qualifications, data/IP responsibilities, brand rules, project change controls, security and confidentiality, and real acceptance criteria. Confirm who is legally authorized to sign on behalf of each college. A ceremonial MOU without commercial obligations or a clear academic approval route should not be treated as a contract or affiliation.

Track receivables and client concentration. A large college contract with deferred collections creates governance and liquidity exposure, not just revenue.

\subsection{65. Phase G3: considering institutional or venture investors}

Prepare a clean cap table and due diligence room. Define the financing's purpose and whether the business genuinely offers venture-style returns. Simulate dilution over several rounds, option pools and exit values. Examine vetoes, board seats, debt restrictions, forced-sale rights and founder service obligations. Negotiate materiality thresholds and preserve room for routine academic and commercial decisions.

If a technology spinout raises money, ensure the academic core's assets, customer data and IP have not been transferred in ways that violate existing shareholder, student or donor rights.

\subsection{66. Phase G4: research governance}

Appoint qualified external scientific review, define research integrity and ethics procedures, separate sponsor agreements and restricted financing, and codify publication, authorship and invention policies. The founder should not be the final scientific arbiter on every research question. Protect student research participation and conflicts where educators invest in student companies.

Scientific freedom becomes credible when the institution can publish results that are commercially inconvenient.

\subsection{67. Phase G5: a separately governed academic institution}

Before a campus or university proposal, obtain a full legal architecture opinion covering sponsor eligibility, governing bodies, academic authority, land, funding, donor restrictions, name/brand, for-profit affiliate transactions and future succession. The sponsor may lawfully retain certain appointment or constitutional rights, but should not assume it can privately extract university surplus or overrule statutory academic bodies.

A conceptual constellation of commercial and mission entities may be sensible; the lawful implementation depends on actual state and national rules, not a generic diagram.

\subsection{68. A practical Bayes control policy}

Reserve founder or board approval for financing instruments, issuance of equity, business-model changes, major property, research-related IP disposal, substantial related-party transactions and commitments endangering core liquidity. Delegate admissions logistics, teaching schedules, ordinary vendor purchases and support decisions to qualified staff within published standards. Give academic reviewers independent authority to uphold credible assessment criteria.

This preserves strategic control while letting the institution become more capable than the founder acting alone.
